% ECONOMICS_PROBLEM: {"id": "L13-555", "title": "Common ownership and competition", "jel_code": "L13", "source_id": "OP-0079"}
% !TeX program = xelatex
\documentclass[11pt,a4paper]{article}
\usepackage[margin=23mm]{geometry}
\usepackage{fontspec}
\setmainfont{Latin Modern Roman}
\usepackage{amsmath,amssymb,mathtools}
\usepackage{xcolor,enumitem,booktabs,longtable,array,xurl,hyperref}
\definecolor{muted}{HTML}{53636C}
\hypersetup{colorlinks=true,urlcolor=blue,linkcolor=blue}
\setlength{\parindent}{0pt}
\setlength{\parskip}{5pt}
\newcommand{\R}{\mathbb R}
\newcommand{\E}{\mathbb E}
\newcommand{\N}{\mathbb N}
\newcommand{\ind}{\mathbf 1}
\newcommand{\Var}{\operatorname{Var}}
\newcommand{\Cov}{\operatorname{Cov}}
\newcommand{\supp}{\operatorname{supp}}
\DeclareMathOperator*{\argmin}{arg\,min}
\DeclareMathOperator*{\argmax}{arg\,max}
\newcommand{\entrymeta}[3]{{\sffamily\small\color{muted}\textbf{\texttt{#1}}\quad|\quad #2\par #3\par}}
\newcommand{\Problemref}[1]{\texttt{#1}}
\newcommand{\JELref}[2]{\texttt{#2} (JEL \texttt{#1})}
\begin{document}
\section*{Common ownership and competition}
\textbf{Problem ID:} \texttt{L13-555}\quad \textbf{JEL:} \texttt{L13}\par
% BEGIN ECONOMICS STATEMENT
\label{problem:OP-0079}
\label{atlas:prob:I9}

\noindent\textbf{Original identifier:} I9 (atlas item 79). \textbf{Classification:} Ownership influence causal identification research program. \textbf{Original tractability label:} Hard (editorial, not a complexity result).

\subsection*{Definitions and assumptions}

Let $b_{hj}\geq0$ be investor $h$'s cash-flow share in firm $j$, and $c_{hj}\geq0$ its effective influence weight in that firm's decision process. These are distinct inputs. An illustrative manager objective is
\[
\sum_k\lambda_{jk}\pi_k,\qquad
\lambda_{jk}=\frac{\sum_hc_{hj}b_{hk}}{\sum_hc_{hj}b_{hj}},
\]
where the denominator is positive and $\lambda_{jj}=1$. This aggregation requires a specified governance mechanism; portfolio overlap alone does not imply it. Let $O$ describe holdings, influence, communications, and compensation. Product-market conduct, demand, costs, and ownership-induced governance changes form a declared equilibrium model with selection $\sigma$. If $Y$ is welfare, specify a finite real-valued evaluator with declared cardinal normalizations and weights on consumers, workers, and owners, accounting for monetary transfers consistently.

\subsection*{Formal research question}

Identify or sharply bound the causal effects of a defined ownership or governance intervention $O_1$ versus $O_0$ on prices, output, quality, and welfare:
\[
\Delta_Y^\sigma=Y(e_\sigma(O_1))-Y(e_\sigma(O_0)).
\]
Determine what observations and exclusions identify the influence weights or a less restrictive mechanism connecting investors to product decisions. Test whether direct governance measures predict conduct changes beyond market structure and cost changes. Compare the competition channel with governance-efficiency effects under explicitly defined counterfactual restrictions, and report ambiguity when the same price patterns are consistent with several mechanisms.

\subsection*{Interpretation and scope}

Institutional investors can have shared cash-flow interests without coordinating or influencing pricing. Conversely, small holdings may carry substantial influence in a dispersed shareholder base. Modified concentration measures derived from ownership models therefore are model-implied proxies and not direct observations of managerial objectives. Holdings respond to firm prospects, and index changes can affect financing, liquidity, or attention as well as ownership overlap; an instrument requires exclusions covering these channels. Product-market concentration and common ownership often move together, demanding variation or restrictions that separate their roles. Direct voting, engagement, and compensation evidence can support a governance link but may also reflect anticipation of product-market changes. Holding demand and technology fixed is useful for a competition mechanism counterfactual, while total policy evaluation must allow the intervention's governance and financing effects. A successful extension would establish a credible chain from specified investor influence to strategic choices and customer outcomes, with uncertainty over the influence mapping, rather than interpret a single ownership--price regression as universal causal evidence.

\subsection*{Source audit and status}

[S1] reports airline evidence; [S2] documents historical model-implied incentives; [S3] is Kennedy, O'Brien, Song, and Waehrer's 2017 working paper presenting economic foundations and empirical analysis. Its working-paper status is retained, and it is not confused with O'Brien--Waehrer's separately titled critique. All three original author-year entries are resolved. Competing interpretations make this an editorial causal and mechanism research program, not a theorem settled or certified open by one cited estimate.

\subsection*{References}

\begin{enumerate}[label={[S\arabic*]},leftmargin=*]

\item José Azar, Martin C. Schmalz, and Isabel Tecu (2018). \emph{Anticompetitive Effects of Common Ownership}. Journal of Finance 73(4), 1513--1565. \url{https://doi.org/10.1111/jofi.12698}\par \textit{Source role:} Airline evidence supporting common-ownership hypothesis. \textit{Locator:} Abstract and article metadata.

\item Matthew Backus, Christopher Conlon, and Michael Sinkinson (2021). \emph{Common Ownership in America: 1980--2017}. American Economic Journal: Microeconomics 13(3), 273--308. \url{https://www.aeaweb.org/articles?id=10.1257/mic.20190389}\par \textit{Source role:} Historical patterns of model-implied ownership incentives. \textit{Locator:} Abstract and article metadata.

\item Pauline Kennedy, Daniel P. O'Brien, Minjae Song, and Keith Waehrer (2017). \emph{The Competitive Effects of Common Ownership: Economic Foundations and Empirical Evidence}. Working paper, July 24, 2017; SSRN 3008331. \url{https://papers.ssrn.com/sol3/papers.cfm?abstract_id=3008331}\par \textit{Source role:} Critical analysis; working-paper status retained. \textit{Locator:} Abstract and article metadata.

\end{enumerate}

\subsection*{Common conventions: Source mathematical conventions}

Unless an entry states otherwise, agent and firm sets are finite. State and action spaces are measurable, strategies and outcome maps are measurable, probability laws are specified, and expectations exist for the targets under discussion. Infinite-horizon discount factors lie in $(0,1)$ and discounted objectives are finite. Norms, tolerances and complexity input sizes must be specified for computational comparisons. Constants or primitives that appear locally are inputs, not empirical facts asserted by this catalogue.

\textbf{Identification.} A statistical model $m\in\mathcal M$ implies an observable law $P_m$ and target $\theta(m)$. At law $P$, its identified set is
\[
\Theta(P;\mathcal M)=\{\theta(m):m\in\mathcal M,\ P_m=P\}.
\]
Point identification holds when this set is a singleton, or equivalently when $P_{m_1}=P_{m_2}$ implies $\theta(m_1)=\theta(m_2)$ for all compatible models. Sharp bounds mean bounds attained, or approached when the boundary is not attained, by compatible models. Writing this set is a definition, not a solution: the substantive task is to establish informative restrictions and derive or estimate the set. An unrestricted-confounding model may give only trivial bounds.

\textbf{Inference.} Identification, estimability and valid uncertainty quantification are separate questions. Fix $\alpha\in(0,1)$. A confidence procedure $C_n$ over a declared law class $\mathcal P$ has asymptotic uniform coverage at level $1-\alpha$ when
\[
\liminf_{n\to\infty}\inf_{P\in\mathcal P}\Pr_P\{\theta(P)\in C_n\}\geq1-\alpha.
\]
For set-identified targets the coverage event must be defined separately, for example coverage of the full identified set. If an entry uses identification-indexed models instead of uniquely indexed laws, the same coverage convention is applied over those models. Pointwise limits do not establish uniform validity, and asymptotic validity does not guarantee finite-sample precision. Sample sizes, dependence structures and weak-identification sequences are part of each application.

\textbf{Counterfactuals.} $Y_i(a)$ is a potential outcome under a specified intervention, including its timing, implementation and versions. It is not $Y_i$ conditional on voluntarily choosing $a$. A full intervention vector permits interference; shorthand scalar treatment notation does not silently impose no interference. Assignment, selection, attrition, anticipation and measurement must be specified. A claim about an instrument's local effect is not automatically a population policy effect.

\textbf{Economic equilibrium.} A counterfactual model includes best responses, constraints, feasibility, market clearing, state transitions and expectations consistent with the stated information. When several equilibria exist, either a declared selector is an input or the target is set-valued. Comparisons across policy regimes must use the stated selection convention. Reduced-form relationships are not presumed invariant after an equilibrium-changing intervention.

\textbf{Optimization and welfare.} Optimization is over the stated feasible set. If attainment is not established, $\sup$ or $\inf$ and $\varepsilon$-optimal choices replace an asserted maximizer or minimizer. For example, with value $V=\sup_{a\in\mathcal A}W(a)<\infty$, an $\varepsilon$-optimal set is $\{a\in\mathcal A:W(a)\geq V-\varepsilon\}$ for $\varepsilon>0$. Benefits, costs, risk and health or environmental outcomes can be aggregated only using declared units and conversion weights. Interpersonal utility weights, the treatment of fiscal transfers and foreign profits, discounting, and the behavioral welfare criterion are explicit inputs. A comparative-static derivative additionally requires existence and appropriate differentiability of the selected counterfactual.

\textbf{Sources.} Local labels [S1], [S2], and so forth restart in every question. Locators identify the relevant abstract, model, section or result at the level actually checked; a bibliographic or abstract-level verification is not represented as inspection of an entire proof. Working papers are distinguished from later publications and changing versions. Source summaries are paraphrased. All URL checks and editorial formulations refer to the audit date stated above.
% END ECONOMICS STATEMENT
\end{document}
